% Reconstructed text from V5 PDF page 22; non-native source
\par\smallskip\noindent\rule{\textwidth}{0.4pt}\par{\sffamily\footnotesize Original V5 page 22 -- recovered text}\par\begin{flushleft}\footnotesize\ttfamily
\par\smallskip
\par\smallskip
\noindent Diagnostics and Serving Infrastructure for MAS 22\par
\par\smallskip
\par\smallskip
\par\smallskip
\noindent agents appear more capable than they are when evaluated on single attempts rather than re-\par
\noindent peated trials. BenchGuard (COLM 2026) audits benchmark quality itself, examining three ex-\par
\noindent isting benchmarks for annotation consistency, task ambiguity, and reproducibility issues.\par
\noindent  Safety-specific benchmarks span three evaluation dimensions. GuardianAgentBench (580\par
\noindent scenarios, 81 tools, 5 failure types across 3 agent frameworks) is the most comprehensive, pro-\par
\noindent viding multi-agent scenarios with structured failure-type annotations. Agent Security Bench\par
\noindent (ASB, ICLR 2025) evaluates security vulnerabilities across 10 scenarios with 400+ tools, find-\par
\noindent ing that existing defense mechanisms are “barely effective” against adversarial attacks (attack\par
\noindent success above 84\%). TAMAS (ACL 2026) tests adversarial robustness of multi-agent systems\par
\noindent with 300 instances covering Byzantine, colluding, and contradicting agent behaviors. OSGuard\par
\noindent (324 examples, 45 tasks) evaluates OS-level action safety with a guardrail architecture. Agent-\par
\noindent SafetyBench (2,000 test cases) and AgentHarm (440 tasks, ICLR 2025) focus on single-agent\par
\noindent safety evaluation.\par
\par\smallskip
\noindent Gap: no cross-layer benchmark. A striking finding from the benchmark analysis is that\par
\noindent none of the 12 Tier-1 benchmarks evaluates the diagnostics-infrastructure interaction:\par
\noindent how diagnostic trace capture affects serving throughput, or how serving-side eviction policies\par
\noindent affect diagnostic fidelity. TraceElephant’s middleware is described as “lightweight” (§3.2) but\par
\noindent provides no quantitative overhead measurements. Who\&When’s attribution is purely offline\par
\noindent post-hoc analysis. This gap is not merely an oversight — it reflects the architectural separation\par
\noindent between diagnostic and serving research communities, and motivates the validation protocol\par
\noindent proposed in Section 7.\par
\par\smallskip
\noindent Cross-cutting analysis. Figure 4 presents a coverage matrix across 10 evaluation dimensions\par
\noindent for all 12 benchmarks. Three systematic gaps are visible:\par
\noindent  1. No benchmark tests diagnostic-infrastructure interaction. The “Diag-infra interac-\par
\noindent  tion” column is entirely empty: 0 of 12 benchmarks measure how diagnostic trace cap-\par
\noindent  ture affects serving throughput, or how eviction policies affect attribution accuracy. This\par
\noindent  is the strongest visual evidence for the survey’s gap thesis.\par
\noindent  2. No benchmark tests concurrent agents. All attribution benchmarks evaluate single-\par
\noindent  agent or sequential multi-agent traces; none evaluates attribution under concurrent agent\par
\noindent  execution where KV cache contention and scheduling decisions may affect trace avail-\par
\noindent  ability.\par
\noindent  3. Attribution and scale are inversely related. Only 4 of 12 benchmarks provide step-\par
\noindent  level attribution labels (TraceElephant, Who\&When, TrajAudit). The largest benchmarks\par
\noindent  (Who\&When Pro at 12,326 traces, Agent-SafetyBench at 2,000 cases) have no attribution\par
\noindent  labels. The annotation cost barrier (TraceElephant’s \$12.66/trace, Who\&When’s 84.3\par
\noindent  human-hours) explains this inverse relationship: high-fidelity attribution at large scale\par
\noindent  would exceed \$150K.\par
\noindent  Trace fidelity is bipartite: benchmarks either provide full execution traces (TraceElephant,\par
\noindent Who\&When, RootSE) or output-only evaluation (ReliabilityBench, Agent-SafetyBench, Agen-\par
\noindent tHarm, OSGuard). The intermediate case — partial trace with structured metadata but without\par
\noindent full step-level tool logs — is absent, despite being the most common deployment scenario.\par
\noindent Failure-attribution benchmarks and safety benchmarks use entirely different evaluation proto-\par
\noindent cols: attribution benchmarks measure causal agent/step identification accuracy, while safety\par
\noindent benchmarks measure attack success rate or scenario coverage. No benchmark evaluates both\par
\noindent dimensions simultaneously.\par
\end{flushleft}
